\subsection{NVOS: evidence construction versus source resolution}
We test a fixed $2\times2$ design on all eight tasks. The old evidence recipe
bilinearly samples logits, thresholds their sigmoid probabilities at 0.5, and erodes positive and
negative interiors. The alternative pools native probability $p$ with confidence
$c=|2p-1|$, using target $\operatorname{area}(cp)/\operatorname{area}(c)$ and
weight $\operatorname{area}(c)$. Zero confidence contributes no evidence.
Both use original tracked negative reliability and geometry visibility.
This changes pooling, boundary authority and confidence weighting together;
it does not isolate one of those factors individually.

Source grids are $256\times192$ and $512\times384$. The old recipe's erosion
radius scales from two to four pixels to retain its normalized image footprint.
Changing the grid changes contribution sampling and observed support, with
geometry parameters fixed. Teachers, tracks, prompts, legal frames, 120 Huber
updates and final threshold 0.5 are shared. A CPU CSR solver first reproduces
the retained low-resolution field and its final IoU. Every task/arm is predicted
before scoring. Other benchmark outputs are unchanged hash-verified references.

\begin{table}[ht]
\centering\small
\begin{tabular}{lrrrr}
\toprule
Task & Old / low & Coverage / low & Old / high & Coverage / high \\
\midrule
fern & 82.19 & 82.63 & 81.87 & 83.53 \\
flower & 86.53 & 86.23 & 86.33 & 87.38 \\
fortress & 97.23 & 96.34 & 97.16 & 97.46 \\
horns\_center & 95.73 & 94.33 & 95.52 & 95.55 \\
horns\_left & 92.70 & 93.01 & 91.71 & 94.43 \\
leaves & 94.08 & 94.06 & 91.01 & 95.06 \\
orchids & 68.14 & 68.09 & 67.68 & 70.05 \\
trex & 75.51 & 78.95 & 73.63 & 78.75 \\
\midrule
Mean (retained GT) & 86.51 & 86.71 & 85.61 & 87.78 \\
Mean (LUDVIG GT) & 86.68 & 86.96 & 85.78 & 87.93 \\
AP (retained GT) & 98.02 & 98.16 & 97.85 & 98.15 \\
AP (LUDVIG GT) & 97.99 & 98.13 & 97.81 & 98.12 \\
\bottomrule
\end{tabular}

\caption{Fixed NVOS factorial study. Each column applies one recipe to all
eight tasks. Task rows use retained GT; mean rows separately report both GT boundary
conventions, followed by AP on the exact same response images. The two horns tasks share one scene. No per-task winners or
GT-selected thresholds are used.}
\label{tab:nvos-coverage}
\end{table}

The cohort has already informed development, so prediction sealing does not
make this an independent final test. Results and source erosion statistics
are retained per task; uncertainty estimates resample seven scene clusters.
These are candidate experiments, not an automatic change to the v10 default.

Coverage at the finer grid raises mean IoU from 86.5134 to 87.7776
(+1.2642 points), improving seven tasks and decreasing horns-center by 0.1731.
Mean AP increases from 98.0222 to 98.1527. Increasing resolution alone decreases
IoU on every task (mean 85.6138); coverage alone at the coarse grid gives 86.7069.
The factorial interaction is +1.9703 points. Thus finer sampling is useful
with the changed evidence construction, not a sufficient explanation by itself.
The seven-scene cluster bootstrap gives a descriptive 95\% interval of
[0.7381, 2.0202] points for the combined change; this is not an independent-test
confidence guarantee after development on this cohort.

Trex improves by 3.2386 IoU and 2.2656 AP points; orchids improves by 1.9049 IoU
but loses 2.0161 AP points. The latter warns that foreground recovery at a fixed
threshold does not ensure improved global ranking. Flower remains at 87.3824 IoU,
so missed foreground extent is not resolved. Using LUDVIG's GT convention gives
87.9280, versus its local three-seed reproduction at 91.2577, with the previously
declared input and geometry differences. No target-RGB benefit is isolated here.

\IfFileExists{tables/nvos_inherited_readouts.tex}{
\paragraph{Recovering the full capability of v10.}
The stored v10 output includes an independently fitted selection in addition
to the continuous response used in its original NVOS main score. We freshly
render both retained outputs for all eight tasks using their existing fixed
readout rules, without fitting a new model. Table~\ref{tab:nvos-inherited}
reports both. The evaluation quantity remains binary-mask IoU on the same GT;
the output operator changes from a full-scene response to a selected-set alpha
projection. Their difference measures existing output capability, not the
gain of the new coverage construction or the optional selector.
The inherited independent set obtains 85.2245 IoU, below the original continuous
response at 86.5134; no higher aggregate was hidden by reporting only that response.
\begin{table}[ht]
\centering\small\begin{tabular}{lrrr}
\toprule
Task & v10 response & v10 independent set & $\Delta$ (pp) \\
\midrule
fern & 82.19 & 81.71 & -0.47 \\
flower & 86.53 & 85.66 & -0.87 \\
fortress & 97.23 & 95.93 & -1.30 \\
horns\_center & 95.73 & 94.06 & -1.67 \\
horns\_left & 92.70 & 90.63 & -2.07 \\
leaves & 94.08 & 93.72 & -0.37 \\
orchids & 68.14 & 64.04 & -4.10 \\
trex & 75.51 & 76.04 & +0.53 \\
\midrule Mean (retained GT) & 86.51 & 85.22 & -1.29 \\
\midrule Mean (LUDVIG GT) & 86.68 & 85.39 & -1.29 \\
\bottomrule
\end{tabular}

\caption{Our two existing v10 NVOS outputs, scored on all eight tasks. One
fixed output rule is used throughout each column. The historical continuous
response score remains preserved; no scene-wise output switching is used.}
\label{tab:nvos-inherited}
\end{table}
}{}

\IfFileExists{tables/nvos_readout_factorial.tex}{
A field-by-readout control further crosses the two stored fields with the two
output operators. All eight geometries and observed-support masks are identical
between the two fields. The continuous operator retains threshold 0.5; the
selected-alpha operator retains its existing membership and alpha thresholds.
This separates the effect of the fitted member field from the composite output
operation (including its fixed readout threshold). All arms are declared and
predicted before scoring, with no newly optimized field.
\begin{table}[ht]
\centering\small\begin{tabular}{llrrrr}
\toprule
Field & Readout & IoU & AP & Precision & Recall \\
\midrule
Semantic & Full scene & 86.51 & 98.02 & 95.70 & 90.10 \\
Semantic & Selected alpha & 84.39 & 95.41 & 87.16 & 96.70 \\
Member & Full scene & 6.91 & 94.38 & 97.96 & 6.95 \\
Member & Selected alpha & 85.22 & 92.77 & 96.20 & 88.24 \\
\bottomrule
\end{tabular}

\caption{Inherited field versus output operator, macro over all eight NVOS tasks
using retained GT. Crossed readouts are diagnostic controls, not task-wise winners.}
\label{tab:nvos-readout-factorial}
\end{table}
With the semantic field fixed, selected-alpha projection raises recall from
90.0997 to 96.6960 but reduces precision from 95.6980 to 87.1647 and IoU to
84.3934. Conversely, the independently fitted member field is not calibrated
as a full-scene semantic response: that crossed readout gives 6.9147 IoU despite
94.3836 AP. Its intended selected-alpha readout recovers 85.2245 IoU.
These controls demonstrate an objective/readout compatibility requirement,
not a failure of the existing member field at its intended task. Restoring
complete output capability requires preserving this pairing.
}{}
